We start our analysis by looking at the questions from the perspective of utilitarianism.  Our analysis is based on~\cite{odihi20}.

At a first glance, profiling does not seem to have an impact on the total utility of the agents at play. Indeed, if the health insurance company at hand is simply redistributing risk among the customers without making a profit (i.e.~the company is only moving money around), then whatever utility is lost by some agent by having to spend more would then be gained by everyone else having to spend a bit less.

The analysis stays the same if a percentage of the incoming revenue is taken by the company as profit or running costs. Of course, this assumes the utility of the customers is a linear function with respect to the amount of money spent. Were it the case this function was concave, then utilitarianism would side against profiling (this is easy to see in the trivial case where all the individuals have equal income). Conversely, utilitarianism would side with profiling in the case of convex utility functions. According to~\cite{prospect_theory}, the utility functions observed in practice vary from individual to individual, thus we will avoid the issue by simply assuming linear functions.

In practice, statistical profiling introduces additional costs — data collection, statistical analysis and the disputes involved introduce additional costs for the insurance companies, costs which must then in turn be distributed to the customers. Thus, if the total amount of money required for the functioning of the system is the only utility source considered for the analysis, then utilitarianism comes out against profiling.

Of course, the total amount of money spent is hardly the only matter worth taking into consideration. For one, people who have to pay more because of circumstances outside their control might feel discriminated against or stigmatized. Moreover, at the individual level, is the exchange of utility equivalent? That is, does a person who get sick but gets reimbursed by the insurance company remain utility-neutral on the matter? Does a person who is not sick, but keeps paying for insurance gain enough utility from the sense of security involved to counteract the loss of utility from having to spend a given amount of money?

Indeed, when looking at a company that profiles its customers, a person that has greater-than-average risk will produce additional negative utility not only when getting sick (this is always the case, be it with or without profiling), but also from potential feelings of discrimination or from feeling stigmatized. If the cause of the increased risk is one the person can influence, then the distress caused by the loss of utility could act as an additional push for the agent at hand to improve their condition, generating positive utility in the long term, although as discussed in the previous sections, many of the risks considered do not fall into this category. On the other hand, in a world without statistical profiling, people who has smaller-than-average risk will produce additional negative utility from the feeling they are paying for other people's health. Weighting the two sources of negative utility against each other without imposing further restrictions on the exact values of the utilities at hand is impossible. We will thus have to move our focus to more specific classes of risk factors.
